\documentclass[../../../include/open-logic-section]{subfiles}

\begin{document}

\olfileid{sth}{spine}{zf}
\olsection{$\Z$ او $\ZF$: يو مهم پړاو}

د بنسټ له اصل سره يو بل مهم پړاو ته رسېږو۔ $\Zminus$ او $\ZFminus$
تيورۍ مو څېړلې وې؛ ويلي مو وو چې دا د کوم ټاکلي شي له ايستلو وروسته
پاتې تيورۍ دي۔ هماغه ټاکلے شوے شے د بنسټ اصل دے۔ نو:

\begin{defn}
$\Z$ تيوري $\Zminus$ ته د بنسټ اصل ورزياتوي۔ نو اصول ئې دا دي: د غړو
له مخې برابري، اتحاد، جوړې، د فرعي سټونو سټونه، نامتناهيت، د بنسټ اصل،
او د بېلونې د شېما ټولې بېلګې۔

$\ZF$ تيوري $\ZFminus$ ته د بنسټ اصل ورزياتوي۔ په بله وينا، $\ZF$
د تعويض ټولې بېلګې $\Z$ ته ورزياتوي۔
\end{defn}

خو ښايي بيا هم يو سوال درته پيدا شوے وي۔ که منظموالی په $\ZFminus$
کښې د بنسټ له اصل سره همارزښته دے او د منظموالي توجيه څرګنده ده، نو
ولې دا اوږده لاره ونيسو او د منظموالي پر ځاے د بنسټ اصل خپل بنسټيز
اصل وګرځوو؟

تاريخي دليلونه، يعنې دا چې چا کومه وينا کله جوړه کړې وه، يوې خوا ته
پرېږدئ۔ اصلي دليل دا دے چې د بنسټ اصل د $V_\alpha$ د تعريف له کارولو
پرته وړاندې کېدلے شي۔ هغه تعريف د \olref[recursion]{sec} پر ټول کار
ولاړ و: د نامتناهيت هاخوا بازګښت ثبوت ته اړ وو څو د تعريف توجيه وښيو۔ خو
د نامتناهيت هاخوا بازګښت په خپل ثبوت کښې مو \emph{تعويض} کارولے و۔ نو
که څه هم د بنسټ اصل او منظموالی د $\ZFminus$ په شتون کښې همارزښته دي،
د $\Zminus$ په شتون کښې همارزښته نۀ دي۔

په حقيقت کښې موضوع تر دې ساده څرګندونې هم جدي ده۔ که څه هم دا بحث
د دې کتاب له حدودو ډېر بهر دے، معلومېږي چې $\Zminus$ او $\Z$ دواړه
د $V_\alpha$ د تعريفولو دپاره ډېرې کمزورې دي۔ نو که يوازې په $\Z$ کښې
کار کوئ، منظموالی په هغې بڼې کښې چې موږ بيان کړے، آن \emph{معنا} نۀ
لري۔ همدا وجه ده چې زموږ رسمي اصل د منظموالي پر ځاے د بنسټ اصل دے۔

له دې وروسته به، که بل ډول مو نۀ وي ويلي، له کومې اضافي يادونې پرته
په $\ZF$ کښې کار کوو۔

\end{document}
